\chapter{Color grading}
\section{Choisir une zone tonale}
Les trois boutons sélectionnent les ombres, les tons moyens et les hautes
lumières. Chaque zone conserve sa teinte et son intensité indépendamment des
deux autres. Contrairement au mélangeur, la sélection porte ici sur une zone
tonale, pas sur les objets déjà d’une certaine couleur.

\section{Utiliser la roue}
Cliquer dans la roue puis faire glisser la sélection. L’angle choisit une
teinte et l’éloignement du centre règle l’intensité. Pendant un glissement,
sortir du bord conserve la manipulation à l’intensité maximale ; il n’est
pas nécessaire de reprendre la poignée à chaque tour.
Le contrôle \emph{Intensité}, de 0 à 100, permet aussi une saisie numérique.
À zéro, la coloration de la zone est neutralisée même si une teinte reste mémorisée.

\paragraph{Exemple.}
Donner une faible teinte bleutée aux ombres et une faible teinte chaude aux
hautes lumières. Vérifier ensuite les tons moyens et les carnations, puis
comparer avec \touche{Espace} maintenu dans la vue. Éviter de pousser les
trois zones au maximum pour juger un rendu naturel.

\paragraph{Clavier.}
Le curseur d’intensité accepte les flèches, la saisie et le double-clic de
réinitialisation décrits pour les réglages. Aucun pilotage angulaire de la
roue par les flèches n’est actuellement déclaré.
